\section{6교시 (75분) — 실습 ② 이슈 로그·품질 기록·감리 대응·검수 보류 회신}

% =====================================================================
% 6교시 — 실습 ②
% =====================================================================
\begin{frame}{6교시 — 이 교시에서 하는 것}
\begin{itemize}
\item 5교시의 세 문서(⑬·⑭·⑮)를 \textbf{손으로} 쓴다 — 워크북 §3\til{}5
\item 세 과제, 한 사건 — 감리 지적 3·4번(IS-07)이 검수 보류·대기 비용·R-20으로 번진다
\item ⑬ 이슈 2건 신규 기입 + 리스크 1건 재평가 (25분)
\item ⑭ 감리 지적 3·4번 원인 귀속 이의 서면 (20분)
\item ⑮ 구두 검수 보류에 서면 회신 + CR-08·11 변경 대가 (25분)
\end{itemize}
\keymsg{세 문서는 같은 사건을 세 자리에서 기록한다 — 귀속 칸 하나가 세 문서의 숫자를 정한다}
\note{4교시 실습 ①과 같은 운영 — 개인 작성 후 조별 대조, 강사는 순회. 다른 점은 세 과제가 한 사건(09-25 감리 1차 지적 8건, 그중 3·4번의 "수행사 설계 미흡" 기재)에서 갈라진다는 것이다. ⑬은 이 사건이 이슈 로그(IS-07)와 등록부(R-20)에 어떻게 들어가는지, ⑭는 귀속을 어떻게 뒤집는지, ⑮는 그 결과가 검수와 대기 비용에 어떻게 이어지는지를 쓴다. 시간이 부족하면 ⑮의 변경 대가 산정표는 두 행(CR-08·CR-11)만 쓰게 한다. 워크북 §3.7·§4.7·§5.7 기본 과제와 다르게, 오늘 교실에서는 완성 예시본이 다루는 사건 자체를 재구성한다(예시본 숫자를 베끼지 않고 근거 문서에서 다시 도출). 예상 소요 3분.}
\end{frame}

\begin{frame}{과제 ⑬ — 워크북 §3: 이슈 2건 신규 기입 + 리스크 1건 재평가}
\begin{itemize}
\item 시점 2026-11-30, G2 직후 — 워크북 §3.7 기본 과제 상황
\item IS-12: 영업본부 "델리 서브 브랜드 상품 코드 체계 변경" 정식 요구 (R-10 실현)
\item IS-13: P2 벤더 MC 2027-03-19로 2주 추가 지연 통보 (R-01 잔여 실현)
\item 재평가: R-10 행 — 확률·영향·EMV·컨틴전시 0.80과 부족분 자금원
\item 항목 6 잔여 EMV 대조 — 7.26 \ra{} ? / 허용범위 8.0 넘는가
\end{itemize}
\keymsg{이슈 행은 열 아홉 개를 다 채운다 — 특히 "출처"(R-nn 전환 / 신규 발견)와 "기한"·"에스컬레이션 단계"}
\note{워크북 §3.3 항목 2(이슈 로그 열 규칙)와 항목 3(리스크 \ra{} 이슈 전환 규칙)을 펴게 한다. IS-12는 R-10 실현이므로 등록부 R-10 행을 "실현 \ra{} IS-12"로 바꾸고, 재발 가능성이 있으면 같은 번호로 잔여 리스크를 유지한다(R-05 선례, §3.5 항목 4). R-10의 3.8억 전액 실현은 컨틴전시 배정 0.80으로 부족하므로 부족분 자금원 순서(미배정 풀 \ra{} 관리예비비)와 CFO 병행 결재 조건(5억 초과 아님, 5,000만 원 이상 내부통제)을 적게 한다. IS-13은 R-01 잔여 실현이자 관리예비비 대상이므로 당일 스폰서·프로그램 PMO장 보고(Day2 §5.5 항목 9)가 에스컬레이션 칸에 들어가야 한다. 잔여 위협 합이 8.0을 넘으면 예외 보고 대안 3개(대응 강화 / 관리예비비 이관 / 허용범위 재협상)와 결정자(RACI 20·21행)까지. 예상 소요 4분(설명) + 25분(작성).}
\end{frame}

\begin{frame}{과제 ⑭ — 워크북 §4: 감리 지적 3·4번 원인 귀속 이의 서면}
\begin{itemize}
\item AF-1-03 (중대): P2 설비 인터페이스 규격 설비측 확인 서명 부재 — 감리 기재 "수행사"
\item AF-1-04 (보통): 가맹 SW 배포 범위(1.1.6) 설계 미확정 — 감리 기재 "수행사"
\item 대응 관리대장 항목 6의 두 행을 채운다 — 발주사 판정 귀속·근거·조치·기한·이의
\item 이의 서면 1장 — 수신 수석감리원, 참조 스폰서·PMO장·수행사 PM
\item 근거는 두 곳뿐 — WBS 담당 열(발주/수행/공동)과 RACI 행
\end{itemize}
\keymsg{지적 사실은 인정하고 귀속만 다툰다 — 등급·내용에 대한 이의가 아님을 첫 문장에 쓴다}
\note{워크북 §4.3 항목 6과 §4.5 항목 6의 대장 양식을 쓴다. AF-1-03의 근거는 규격이 WBS 1.3.1 발주 담당 산출물이고 P2 전달 책임이 RACI 9행 프로그램 PMO장(A)이라는 것, 그리고 설비 벤더의 05-29 변경 요청에 대한 PMO 회신이 미결이라는 사실 기록(회신 요청 06-05·07-10·08-21, 수행사 서면 통지 06-08·08-24). AF-1-04는 1.1.6 대응안 결정이 RACI 11행 오정근(A)이며 05-08 회신으로 결정을 유보했다는 기록. 서면의 네 단락 — 사실 인정 / 귀속 재분류 요청과 근거 / 발주사 조치 계획(10-16 규격 회신, 11-27 대응안 결정) / 요청 사항(정본 정정, 10-08 회신, 향후 귀속 판정 기준 합의). 교재 §4.3 "왜 귀속 두 글자가 등급보다 무거운가"를 상기시킨다. 예상 소요 3분(설명) + 20분(작성).}
\end{frame}

\begin{frame}{과제 ⑮ — 워크북 §5: 구두 검수 보류에 서면 회신 + 변경 대가}
\begin{itemize}
\item 사건 — 09-28 주간회의 구두 "감리 지적 정리 후 검수" \ra{} 수행사 09-29 서면 "발주사 사유 보류, 대기 비용"
\item 회신 서면 1장(항목 4) — 검수는 §14대로 진행 중, 보류가 아님을 먼저
\item 발주사 사유 기간을 발주사가 스스로 적는다 — 09-28\til{}10-08, 8영업일, 영향 범위 R2 FP 12\%
\item 변경 대가 산정표(항목 2) — CR-08 +90 FP / CR-11 $-$60 FP, 단가 55만 원
\item 항목 6 귀책 분해표에 한 행 — "AF-1-03·04 귀속 정리 8일, 발주사 단독, 병행 가능"
\end{itemize}
\keymsg{구두에는 서면으로, 보류에는 조건으로, 대기 비용에는 기간·범위·판정 절차로 답한다}
\note{워크북 §5.3 항목 2·4·6과 §5.5 완성본을 쓴다. 회신의 뼈대 — ① 검수 요청 접수 사실과 §14 기산일(09-25)·통보 기한(10-16) ② 구두 통보의 정정(보류가 아님) ③ 검수 완료 조건(수행사 귀속 3건 조치) ④ 발주사 사유 기간의 명시(8영업일, 관련 산출물 한정 약 12\%, 나머지와 R3 병행 가능) ⑤ 대기 기록 10-09 제출 \ra{} 10-14 변경협의체 판정 ⑥ 요청 사항. 변경 대가는 CR-08 +90 × 55만 = +0.495억(R4 검수 후 지급), CR-11 $-$60 × 55만 = $-$0.33억(감액, 잔금 차감) — 감액도 합의서 6호에 같은 행으로. 시간이 남는 조는 CR-11의 데이터 전환 용역 $-$2.00(DT-01)까지 항목 7 계약 변경 로그에 넣게 한다. 예상 소요 4분(설명) + 25분(작성).}
\end{frame}

\begin{frame}{작성 요령 ① — 원인 귀속 칸이 책임 귀속을 정한다}
\fig[0.58]{../그림/PM_Day3/fig_4_3_audit_attribution.pdf}
\figcap{그림 4-3. 감리 지적 8건 — 감리 기재 귀속(수행사 6·발주사 2) vs 발주사 판정 귀속(수행사 3·발주사 4·공동 1). AF-1-03·04의 이의 근거는 WBS 담당 열과 RACI 행 두 가지뿐이다}
\keymsg{감리 기재를 그대로 받은 대장은 검수·변경협의체에서 발주사의 미결정을 수행사의 미흡으로 기록한다}
\note{교재 §4.3. 대장에 "감리 기재 귀속"과 "발주사 판정 귀속" 두 열을 따로 두는 이유를 그림으로 보인다 — 한 열이면 감리 기재가 곧 판정이 된다. 판정 근거 열에는 WBS 사전 카드의 담당(발주/수행/공동)과 RACI 행 번호, 그리고 일자가 붙은 사실 기록(발송·회신 요청·서면 통지)만 적는다. 형용사("명백히", "당연히")는 근거가 아니다. 8건 중 이의는 2건뿐이고 나머지 6건은 그대로 받는다 — 전부 다투면 신뢰를 잃는다. 이 두 열의 차이가 ⑮의 귀책 분해표(발주사 단독 8일)와 R-20의 EMV로 이어진다는 것을 짚는다. 예상 소요 4분.}
\end{frame}

\begin{frame}{작성 요령 ② — 서면 회신에 반드시 들어갈 것}
\begin{tbl}[\scriptsize]{L{3.2cm}L{5.8cm}L{2.6cm}}
항목 & 온담 R2 회신(10-01)에서 & 없으면 \\
\midrule
접수 사실·계약 조항·기한 & 09-25 접수, §14 15영업일, 통보 기한 10-16 & 기산일 분쟁 \\
구두 통보의 정정 & "보류"가 아니라 "진행 중, 조건부" & 수행사 기록만 남음 \\
완료 조건 & 수행사 귀속 3건 조치(10-16\til{}10-30) & 검수 무기 연기 \\
발주사 사유 기간 & 09-28\til{}10-08 8영업일, 범위 R2 FP 12\% & 전 기간·전 인력 청구 \\
판정 절차·시한 & 대기 기록 10-09 \ra{} 변경협의체 10-14 & 협상 근거 없음 \\
\end{tbl}
\keymsg{발주사가 스스로 "8영업일"을 적는 이유 — 기간과 범위를 먼저 정의하는 쪽이 협상의 자를 쥔다}
\note{교재 §5.5 "회신의 구조"와 "왜 발주사가 스스로 8영업일을 적는가". 수강생이 가장 저항하는 부분이 넷째 행이다 — 발주사 사유를 왜 발주사가 먼저 인정하느냐. 답은 그림 5-5의 기록 A(수행사만 기록)와 기록 B(발주사 회신)의 차이다. 기록 A만 있으면 대기 기간은 수행사가 정한 기간 전체(09-28부터 무기한)와 전 인력이 되고, 기록 B가 있으면 8영업일 × 관련 산출물 12\%로 한정된다. 인정하지 않는 것이 아니라 한정하는 것이다. 참조 수신자(구매팀장·재경팀장·수석감리원)를 넣는 이유도 같다 — 회신이 계약·지급·감리 세 기록에 동시에 들어간다. 예상 소요 4분.}
\end{frame}

\begin{frame}{작성 요령 ③ — 변경 대가 = FP 델타 × 55만 원, 감액도 같은 표에}
\begin{tbl}[\scriptsize]{L{3.0cm}rrL{1.6cm}L{2.2cm}L{2.0cm}}
CR / 건 & FP $\Delta$ & 금액(억) & 증/감 & 합의·서면 & 지급 시점 \\
\midrule
CR-08 파일럿 임시 연동 & +90 & +0.495 & 증액 & 09-30 · 6호 & R4 검수 후 \\
CR-11 5년 초과 이력 제외 & $-$60 & $-$0.33 & \textbf{감액} & 09-30 · 6호 & 잔금 차감 \\
데이터 전환 용역(CR-11) & — & $-$2.00 & 감액 & 09-25 · DT-01 & 11.4 \ra{} 9.4 \\
\end{tbl}
\begin{itemize}
\item 네 가지가 다 있어야 한 행 — FP 산정 근거 · 단가 · 변경협의체 합의 일자 · 서면 번호
\item FP 확인 3인 서명(PMO 기준선 담당·수행사 PM·형상 담당) — CR-07\til{}11은 BL-02 확인서에서 일괄
\item 대기 비용은 FP가 아니라 인일 × 부속서 단가 — 별도 열, 같은 표
\end{itemize}
\keymsg{감액을 변경으로 안 쓰면 G4 정산에서 수행사가 "합의한 적 없다"고 말할 수 있다}
\note{교재 §3.5와 워크북 §5.3 항목 2. 고정가 FP 계약에서 변경 대가는 산식이 하나뿐이라 쉬워 보이지만, 넷 중 하나가 빠진 행이 G4에서 다시 협상된다(§5.7 점검 질문 1). CR-11은 SI 계약에서 $-$60 FP 감액이면서 데이터 전환 용역에서 $-$2.00 감액이라 두 계약·두 서면(SI 합의서 6호·DT-01)에 각각 들어간다 — 한 변경이 두 행이다. 대기 비용(R-05 0.42, 이번 AF-1-03·04 건은 10-14 판정 후)은 FP 델타가 아니므로 "대기 비용" 열을 따로 두되 같은 표에 쓴다 — 그래야 누적(SI +1.4375 = FP +185 · 대기 0.42)이 한 줄로 나온다. 예상 소요 3분.}
\end{frame}

\begin{frame}{흔한 오류 — 실습 ②에서 반복되는 네 가지}
\begin{itemize}
\item \textbf{이슈를 리스크 등록부에 넣는다} — 실현된 것은 확률 100\%가 아니라 이슈, 등록부 행은 "실현 \ra{} IS-nn"
\item \textbf{귀속 칸을 비운다} — "추후 협의"는 감리 기재가 판정이 되는 것과 같다
\item \textbf{구두 통보에 구두로 답한다} — 수행사 서면(09-29)만 남고 발주사 기록은 0건
\item \textbf{감액을 변경으로 안 쓴다} — CR-11 $-$60 FP가 로그에 없으면 잔금 차감 근거가 없다
\item 기한 없는 이슈 — 이슈가 아니라 배경이다 (워크북 §3.7 점검 질문 3)
\end{itemize}
\keymsg{네 오류의 공통점 — 기록해야 할 순간에 기록하지 않아 상대방의 기록이 사실이 된다}
\note{순회 중 발견한 오류를 이 프레임에서 모아 말한다. 첫째는 §4.4 전환 규칙 — 트리거 발동은 신호(등록부에 관측 사실 기록), 실현은 사건(이슈 로그 등재 + 컨틴전시 집행 개시). 둘째는 귀속 칸을 "공동"으로 얼버무리는 변형도 포함한다 — AF-1-01처럼 근거가 있어 공동이면 되지만 근거 없는 공동은 빈칸이다. 셋째는 mum effect(§5.2)의 계약판 — 나쁜 기록을 남기기 싫어 서면을 피하면 상대의 서면이 유일한 기록이 된다. 넷째는 CR-11이 "좋은 소식"이라 로그를 건너뛰는 경우 — 감액도 변경이고(§3.5), 서면 없는 감액은 G4에서 사라진다. 다섯째는 덤 — 기한과 다음 에스컬레이션 단계가 없는 이슈 행은 완성이 아니다. 예상 소요 4분.}
\end{frame}

\begin{frame}{예시본 참조 — 워크북 완성 예시본과 템플릿 PDF}
\begin{tbl}[\scriptsize]{L{1.2cm}L{3.4cm}L{3.2cm}L{3.6cm}}
산출물 & 워크북 & 빈 템플릿 PDF & 완성 예시 PDF \\
\midrule
⑬ & §3.5 (이슈 11건·재평가·EMV 7.26) & 13\_이슈리스크갱신\_빈템플릿 & 13\_이슈리스크갱신\_완성예시 \\
⑭ & §4.5 (대장 8건·이의 서면 09-29) & 14\_품질통제기록\_빈템플릿 & 14\_품질통제기록\_완성예시 \\
⑮ & §5.5 (산정표·회신 10-01·귀책표) & 15\_조달계약변경\_빈템플릿 & 15\_조달계약변경\_완성예시 \\
통합 & — & 00\_Day3\_템플릿팩\_빈양식 (46쪽) & 00\_Day3\_템플릿팩\_완성예시 (42쪽) \\
\end{tbl}
\begin{itemize}
\item 경로 — \texttt{교재/템플릿/PM\_Day3/}, 교재·워크북과 다르면 \textbf{워크북 완성 예시본이 정본}
\item 예시본은 대조용이다 — 작성 후에 편다. 먼저 펴면 베끼게 된다
\end{itemize}
\keymsg{예시본과 내 답이 다른 자리에 배움이 있다 — 다른 이유를 한 줄로 적어 두라}
\note{작성 시간이 끝난 뒤(약 70분 시점) 예시본을 펴게 한다. 대조 순서 — ⑬은 §3.5 항목 2의 IS-07 행과 항목 5의 R-20 행(오늘 사건이 어떻게 두 문서에 들어갔는가), 항목 6의 변동 분해(6.78 $-$ 1.98 + 3.13 $-$ 0.67 = 7.26)로 자기 재평가의 산식을 검산. ⑭는 §4.5 항목 6의 AF-1-03·04 두 행과 09-29 서면의 네 단락 — 자기 서면에 "지적 사실의 인정"이 첫 문장에 있는가. ⑮는 §5.5 항목 4 회신의 4항(발주사 사유 기간)과 항목 6 귀책 분해표의 "발주사 단독 8일, 병행 가능, 임계경로 영향 0". 가로 페이지 표(이슈 로그·리스크 재평가·감리 대응 대장·변경 대가·귀책 분해)는 템플릿팩 인쇄본을 쓴다. 예상 소요 5분(대조).}
\end{frame}

\begin{frame}{시간 배분과 심화 과제}
\begin{tbl}[\scriptsize]{L{1.4cm}L{6.2cm}L{2.6cm}}
분 & 활동 & 산출 \\
\midrule
0\til{}5 & 과제 설명 (프레임 2\til{}4) & — \\
5\til{}30 & ⑬ IS-12·IS-13 기입, R-10 재평가, EMV 대조 & §3 항목 2·4·6·7 \\
30\til{}50 & ⑭ 대장 2행 + 이의 서면 1장 & §4 항목 6 \\
50\til{}70 & ⑮ 회신 서면 1장 + 변경 대가 2행 + 귀책 1행 & §5 항목 2·4·6 \\
70\til{}75 & 예시본 대조·정리 & 다른 자리 한 줄 \\
\end{tbl}
\begin{itemize}
\item \textbf{심화 과제(3\til{}7년차)} — 자기 회사 이슈 20건의 출처(등록부에 있었나)와 소유자가 PM인 비율을 세라 — 워크북 §3.7
\item 계약서에서 검수·지체상금·변경 세 조항의 위치를 한 장으로 — 워크북 §5.7
\end{itemize}
\keymsg{75분에 세 문서 여섯 조각 — 완성보다 빈칸이 어디인지 아는 것이 오늘의 산출이다}
\note{시간이 밀리면 ⑮의 변경 대가 두 행을 숙제로 돌리고 회신 서면을 지킨다 — 서면 회신이 이 교시의 핵심 산출이기 때문이다. 조별 대조는 4교시처럼 옆 조와 교환해 ⑭ 서면을 상대 조가 "수행사 PM" 역할로 읽고 반박하게 하면 귀속 근거의 약한 자리가 드러난다(3분). 심화 과제 둘은 워크북 §3.7·§5.7의 것을 그대로 — 등록부에 있었던 이슈 비율 30\% 미만이면 등록부가 위험을 못 보는 것, 90\% 이상이면 이슈 로그가 등록부 복사본. 계약 세 조항 요약은 Day4 종료·정산과 연결된다. 예상 소요 2분.}
\end{frame}

\begin{frame}{6교시 핵심 정리}
\begin{itemize}
\item 실현된 리스크는 이슈다 — 등록부 행은 "실현 \ra{} IS-nn", 재발 가능성만 같은 번호로 남긴다
\item 잔여 EMV는 분해해서 본다 — 종결 $-$ / 신규 + / 재평가 ± (6.78 \ra{} 7.26, 여유 0.74)
\item 감리 대장의 "발주사 판정 귀속" 열 — 근거는 WBS 담당 열과 RACI 행, 사실은 일자와 함께
\item 구두 검수 보류에는 서면으로 — 조항·기한·조건·발주사 사유 기간·판정 절차 다섯 가지
\item 변경 대가 한 행 = FP $\Delta$ × 55만 + 합의 일자 + 서면 번호, 감액도 같은 표
\end{itemize}
\keymsg{통제의 마지막 층은 기록이다 — 누가 언제 무엇을 결정하지 않았는지를 발주사가 먼저 적는다}
\note{Day3 산출물 5종이 모두 결정 요청으로 끝난다는 도입의 한 문장으로 돌아간다. ⑬의 결정 요청은 컨틴전시 재배정과 관리예비비 이관(R-10 부족분·R-14·R-17), ⑭는 귀속 정정 회신 10-08과 발주사 조치 2건(10-16·11-27), ⑮는 변경협의체 10-14의 대기 비용 판정과 합의서 6호 서명. 셋 다 결정자가 발주사 계층(프로그램 PMO장·오정근·스폰서)이라는 점 — 수행사에 요구하기 전에 발주사 안의 미결정을 기록하는 것이 Day3 오후의 주제였다. 마무리 30분(제6장 학술 배경·읽을 자료·Day4 예고)으로 넘어간다. 예상 소요 3분.}
\end{frame}
